\documentclass[aspectratio=169,usenames,dvipsnames]{beamer}
\beamertemplatenavigationsymbolsempty
% \documentclass{beamer}

% UNCOMMENT FOR HANDOUTS
% \usepackage{pgfpages}
% \pgfpagesuselayout{2 on 1}[landscape,letterpaper,border shrink=5mm]
% \setbeameroption{show notes}

\mode<presentation>
\usetheme{Boadilla}
\usecolortheme{spruce}
\usepackage{../LizStyle,ulem}
\usepackage{color}
\usepackage[color,matrix,arrow]{xy}
% \input xy
% \xyoption{all}
\usepackage{tikz}
\usepackage{tikz-cd}
\tikzset{commutative diagrams/diagrams={ampersand replacement=\&}}

\AtBeginSection{\frame{\sectionpage}}


% ==========Command for putting comments and removing them=======
% Visible:
\newcommand{\InstructorNotes}[1]{\textit{\color{orange}#1}}
\newcommand{\InstructorList}[1]{
\begin{itemize}
	\color{orange}
	#1
\end{itemize}
}
% Removed:
% \renewcommand{\InstructorNotes}[1]{}
% \renewcommand{\InstructorList}[1]{}


\newcommand{\bvfill}{\vskip0pt plus 1filll}

\usepackage{multimedia}
\usepackage{graphbox}

% \usepackage{algpseudocode}


\graphicspath{ {Figures/} {../Figures/}{../../Figures/} }



\begin{document}
%\pdfoptionpdfminorversion 5
%\logo{\includegraphics[height=1.5cm]{dukeshield.pdf}}
\title[]{Ch 10.3: Convolutional Neural Nets}
\subtitle{Lecture 31 - CMSE 381}
\date{Fri, Nov 22, 2024}



\institute[MSU-CMSE]{Michigan State University\\
			::\\
          Dept of Computational Mathematics, Science \& Engineering}
\author[Dr.~Munch]{Prof.~Elizabeth Munch}

\begin{frame}
 \titlepage


\end{frame}
%--------------------------
\begin{frame}[c]{Announcements }

\begin{columns}
\column{.35\textwidth}
\centering

	\textbf{Last time:}
	\begin{itemize}
	\item Multilayer  NN
	\item pyTorch
	\end{itemize}

	\vspace{.2in}

	\textbf{This lecture:}
	\begin{itemize}
		\item CNNs
\end{itemize}

\textbf{Announcements:}
	\begin{itemize}
		\item Project due Friday 
	% \item Google Colab for last lecture notebook 
	\item HW \#9 is posted
	\item Exam 3 is Dec 4th
	\item Project is due Dec 6th
	% \item Involves running pytorch, so start early!
	
	\end{itemize}

\column{.35\textwidth}
\centering
\includegraphics[width = \textwidth]{Day01-Schedule-Test3.png}

\end{columns}

\end{frame}
%--- Next Frame ---%


%==========
\section{Last time: Neural Nets}
%==========


\begin{frame}[t]{MNIST}

\begin{columns}
\column{.45\textwidth}
\centering
\includegraphics[width = \textwidth, align = c]{Fig10-3}
\column{.45\textwidth}
\centering
\begin{itemize}
  \item Goal: Build a model to classify images into their correct digit class
	\item Each image has $p=28 \cdot 28 = 784$ pixels
	\item Each pixel is grayscale value in [0,255]
	\item Data converted into column order
	\item Output represented by one-hot vector $Y = (Y_0,Y_1,\cdots,Y_9)$
	\item 60K training images, 10K test images
\end{itemize}

\end{columns}

\end{frame}
%--- Next Frame ---%
\begin{frame}[t]{Neural network architecture for MNIST}

\begin{columns}
\column{.3\textwidth}
\centering
\includegraphics[width = .6\textwidth, align = c]{Fig10-4}
\includegraphics[width = \textwidth, align = c]{Eq10-10}
\includegraphics[width = \textwidth, align = c]{Eq10-11}
\includegraphics[width = \textwidth, align = c]{Eq10-12}
\includegraphics[width = \textwidth, align = c]{Eq10-13}
\column{.45\textwidth}
\centering
\begin{itemize}
  \item Two hidden layers. 
  \item Softmax for classification output 
  \item We used $L_1$ has 128 units; $L_2$ has 64
	\item 10 output variables due to class labeling
	\item Result is we are training approx 110K weights
\end{itemize}

\end{columns}

\end{frame}
%--- Next Frame ---%
\begin{frame}[t]{MNIST learning}
\begin{columns}
\column{.2\textwidth}
\centering
\includegraphics[align = c, width = \textwidth, trim = 0 0 220 0, clip]{MNIST-ExamplePrediction.png}

\column{.5\textwidth}
\centering
\includegraphics[width = \textwidth, align = c]{Fig10-4}
\column{.2\textwidth}
\centering
\includegraphics[align = c, width = \textwidth, trim = 210 0 0 0, clip]{MNIST-ExamplePrediction.png}
\end{columns}
\end{frame}
%--- Next Frame ---%

%==========
\section{Convolutional Neural Network }
%==========

\begin{frame}[t]{Last time: Flattening the image}
\begin{columns}
\column{.45\textwidth}
\centering

\[
\begin{pmatrix}
	1 & 1 & 0 \\
	4 & 2 & 1 \\
	0 & 2 & 1
\end{pmatrix}	
\longrightarrow
\begin{pmatrix}
	1 \\ 1 \\ 0 \\
	4 \\ 2 \\ 1 \\
	0 \\ 2 \\ 1
\end{pmatrix}
\]


\column{.45\textwidth}
\centering

\InstructorList{
  \item Just ignore the fact that we have a picture 
  \item Loses a lot of the structure that we might want for more complicated
  \item Answer is CNNs, which take in image data and build a neural net that specifically tracks that structure
}

\end{columns}
\end{frame}
%--- Next Frame ---%

\begin{frame}[t]{Example data set: CIFAR100 Data}

	\centering
\includegraphics[width = .6\textwidth, align = c]{Fig10-5}
\begin{columns}
\column{.45\textwidth}
\centering
\begin{itemize}
	\item 60,000 images: 50K training, 10K  test
	  \item Labels with 20 super classes (e.g.~aquatic mammals)
	  \item 5 classes per super class (beaver, dolphin, otter, seal, whale)
	  \item Images are 32x32
  \end{itemize}

\column{.45\textwidth}
\centering
\footnotesize


\end{columns}

\end{frame}
%--- Next Frame ---%
\begin{frame}[t]{Image channel data}

\begin{columns}
\column{.45\textwidth}
\centering

\InstructorList{
	\item 3 eight bit numbers per pixel (red, green, blue)
	\item Numbers are organized in a \textit{feature map} , which is 3 dimensional array
	\item First two directions are spatial (32x32)
	\item Third is the \textit{channel} axis representing teh three colors
}
\column{.45\textwidth}
\centering
\includegraphics[width = .5\textwidth]{Fig10-8-onlychannels.png}
\end{columns}

\end{frame}
%--- Next Frame ---%
\begin{frame}[t]{CNNs}

\begin{columns}
\column{.45\textwidth}
\centering
\InstructorList{
  \item Idea is to build a neural net that works on this data
	\item Identify low level features on the input image such as small edges, patches of color, etc
	\item Low level features combined into higher-level features (parts of ears, eyes, etc)
	\item Does this using special types of hidden layers: \textit{convolution} layers and \textit{pooling layers}
}

\column{.45\textwidth}
\centering
\includegraphics[width = \textwidth, align = c]{Fig10-6}

\end{columns}

\end{frame}
%--- Next Frame ---%

\begin{frame}[t]{Convolution layer }
	{Convolution Filter}

\begin{columns}
\column{.45\textwidth}
\centering
Original Image:
\begin{equation*}
	\begin{bmatrix}
		a & b & c\\
		d & e & f \\
		g & h & i\\
		j & k & l
	\end{bmatrix}
\end{equation*}

Convolution filter:
\begin{equation*}
	\begin{bmatrix}
		\alpha & \beta \\
		\gamma & \delta
	\end{bmatrix}
\end{equation*}
\column{.45\textwidth}
\centering
Convolved Image
\includegraphics[width = \textwidth, align = c]{CNN-ConvolvedImage.png}
\end{columns}

\end{frame}
%--- Next Frame ---%
\begin{frame}[t]{Convolution Filter Example}

\begin{columns}
\column{.25\textwidth}
\centering
\includegraphics[width = \textwidth, align = c]{CNN-examplematrix}
% \includegraphics[width = .5\textwidth, align = c]{CNN-examplematrix-Mask}

Filter:
\begin{equation*}
	\begin{bmatrix}
		0 & 1 \\
		0 & 1
	\end{bmatrix}
\end{equation*}
\column{.45\textwidth}
\centering

\InstructorList{
	\item DOUBLE CHECK, SHAPE ISN'T RIGHT!
\item $\begin{bmatrix}
	1 & 3 \\
	0 & 3
\end{bmatrix}$
\item Convolved image highlights regions that resemble the convolution filter.
\item Note that the result is smaller than original
}
\end{columns}
\end{frame}
%--- Next Frame ---%
{
\setbeamercolor{background canvas}{bg=Apricot!25}
\begin{frame}[t]{Same example, different filter}
	What is the convolved image?
	\begin{columns}
		\column{.25\textwidth}
		\centering
		\includegraphics[width = \textwidth, align = c]{CNN-examplematrix}
		% \includegraphics[width = .5\textwidth, align = c]{CNN-examplematrix-Mask}
		
		Filter:
		\begin{equation*}
			\begin{bmatrix}
				1 & 1 \\
				0 & 0
			\end{bmatrix}
		\end{equation*}
		\column{.45\textwidth}
		\centering
		
		\InstructorList{
		\item $\begin{bmatrix}
			4 & 3 \\
			3 & 2
		\end{bmatrix}$
		}
		\end{columns}
\end{frame}
}
%--- Next Frame ---%
\begin{frame}[t]{Convolution filter: Bigger example}

	\centering
\begin{columns}
\column{.6\textwidth}
\centering
	\includegraphics[width = \textwidth, align = c]{Fig10-7}

\column{.35\textwidth}
\centering
\InstructorList{
  \item 192x179 image of tiger
	\item Convolution filters are 15x15 images with mostly zero (black) narrow strip of ones (white)
	\item Highlights areas of the tiger that resemble the filter
}

\end{columns}

\end{frame}
%--- Next Frame ---%
\begin{frame}[t]{Convolution layer}
\begin{columns}
\column{.45\textwidth}
\centering
\includegraphics[width = \textwidth, align = c, trim = 0 0 890 0, clip]{Fig10-8}
\column{.45\textwidth}
\centering
\InstructorList{
  \item Image is in color, so 3 channels
	\item This means the convolution filter will have three channels, one per color, each of dimension 3x3
	\item Results are summed to form a 2-dimensional output. Color is "forgotton" at this point
	\item $K$ different filters here means $K$ 2-d output 
	maps, which can then be treated as channels
	\item Padding to make the updated image the same size
}

\end{columns}
\end{frame}
%--- Next Frame ---%
\begin{frame}[t]{More notes on convolution}

\begin{columns}
\column{.45\textwidth}
\centering
\InstructorList{
  \item Filters idea from image processing
	\item Novelty in DL is that the choice of filters are learned
	\item Filter weights: one hidden unit per pixel in convolved image; but more constrained than that
	\item 	
	\item Applying ReLU to convolved image. Sometimes viewed as separate layer in CNN, and then called a detector layer.
}

\column{.45\textwidth}
\centering


\end{columns}

\end{frame}
%--- Next Frame ---%

\begin{frame}[t]{Pooling layers}

\begin{columns}
\column{.45\textwidth}
\centering
\InstructorList{
  \item Condense a large image into a smaller summary image
	\item Lots of options, but max pooling is common:
	\item For each 2x2 non-overlapping block of pixels, get a single new pixel with the maximum value from the block
	\item Provides location invariance: so long as large value in one of the four pixels, the whole block registers as a large value
}

\column{.45\textwidth}
\centering
\includegraphics[width = \textwidth, align = c]{Eq10-3-2.png}
\end{columns}

\end{frame}
%--- Next Frame ---%

\begin{frame}[t]{Putting it together to make a CNN}

\includegraphics[width = \textwidth, align = c]{Fig10-8}
\begin{columns}
\column{.45\textwidth}
\centering
\footnotesize
\InstructorList{
  \item Start with three channels
	\item get a new channel at first hidden layer for each filter. Padding to keep the same size
	\item Next is max pool layer,
	\item convolve \& pool several more times
	% \item Can repeat several convolves before pool, this effectively increases the dimension fo the feature
}

\column{.45\textwidth}
\centering
\footnotesize
\InstructorList{
	\item Since size of image decreased, we increase number of filters in the next convolve layer to compensate
  \item Flattened, then fed into 1 or more fully-connected layers
	\item Last is softmax activation for the 100 classes
}

\end{columns}
\bvfill
\url{https://poloclub.github.io/cnn-explainer/}

\end{frame}
%--- Next Frame ---%

{
\setbeamercolor{background canvas}{bg=Apricot!25}
\begin{frame}[t]{Coding}
\begin{columns}
\column{.45\textwidth}
\centering
\column{.45\textwidth}
\centering
\end{columns}
\end{frame}
}
%--- Next Frame ---%


% \begin{frame}[t]{Data Augmentation}

% 	\includegraphics[width = \textwidth]{Fig10-9}
% \begin{columns}
% \column{.45\textwidth}
% \centering

% \InstructorList{
%   \item Replicate training image many times
% 	\item each time randomly distorted in a natural way that doesn't affect human recognition
% }

% \column{.45\textwidth}
% \centering

% \end{columns}

% \end{frame}
% %--- Next Frame ---%
% \begin{frame}[t]{Results on pretrained classifier: \texttt{resnet50}}

% \begin{columns}
% \column{.45\textwidth}
% \centering
% \includegraphics[width = \textwidth]{Fig10-10}
% \column{.45\textwidth}
% \centering

% \InstructorList{
%   \item Trained on \textit{imagenet} data set
% 	\item Millions of images
% 	\item At left, shows top three categories for the images
% 	\item Note when hawk is zoomed out, labels as fountain
% 	\item Cape weaver is classified as a different tropical bird, jacamar
% 	\item One tool: use some pretrained layers, and then train the last on your particular data
% }


% \end{columns}

% \end{frame}
%--- Next Frame ---%


\begin{frame}[t]{TL;DR}

	\begin{columns}
	\column{.45\textwidth}
	\centering
	\textbf{Feed Forward Neural Net}
	\includegraphics[width = .5\textwidth, align = c]{Fig10-4}
	\includegraphics[width = \textwidth, align = c]{Eq10-2}
	
	\begin{itemize}
		\item Combines input data using learned weights
		\item Linear combo of those to get output
		\item Sometimes softmax to get probability of classification
	\end{itemize}
	
	\column{.45\textwidth}
	\centering
	\textbf{CNN}
	\includegraphics[width = \textwidth, align = c]{Fig10-8}
	
	\begin{itemize}
		\item Specialized NN
		\item Gets next layer via
		\begin{itemize}
			\item Convolution layer
			\item Pooling Layer
			\item Fully connected layer
		\end{itemize}
	\end{itemize}
	
	\end{columns}
	
	
	\end{frame}
	%--- Next Frame ---%

\begin{frame}[t]{Next time}


	

	\includegraphics[align = c, width = .35\textwidth ]{Day01-Schedule-Test3.png}



\end{frame}
%--- Next Frame ---%


\end{document}
